\subsection{属 \texorpdfstring{$49$}{49} 的 \texorpdfstring{$S_4$}{S4} 正则闭包 Jacobian 的商曲线 atlas、Prym 梯与低亏格压缩}\label{subsec:conclusion-s4-quotient-atlas-prym-lowgenus-compression}
沿用
\[
X\longrightarrow \PP^1_u
\]
为文稿既定的属 \(49\) 的 \(S_4\)-正则闭包曲线，
\[
J:=\Jac(X),
\qquad
H^0(X,\Omega_X^1)\cong 5\cdot \varepsilon\oplus 4\cdot V_2\oplus 3\cdot V_3\oplus 9\cdot V_3',
\]
并记 \(A_{\mathrm{sgn}}:=A_{\varepsilon}\)。
在全纯微分的等型分解、标准子群商亏格表以及 \(C_4\triangleleft D_8\) 的九维 Prym 体都已建立之后，可把整条 \(S_4\)-正则闭包的 Jacobian 进一步压缩为一套由根块、商曲线 atlas 与低亏格 Prym 接口共同控制的全局公式。

\begin{theorem}[\texorpdfstring{$S_4$}{S4}--等型 Jacobian 的根块化]\label{thm:conclusion-s4-jacobian-root-block-realization}
设
\[
J\sim_{\QQ}A_{\mathrm{sgn}}\times A_2\times A_3\times A_{3'}
\]
为 \(S_4\)-等型 Jacobian 分解。则存在定义于 \(\QQ\) 上的阿贝尔簇 \(B_2,B_3,B_{3'}\)，满足
\[
\dim B_2=4,\qquad \dim B_3=3,\qquad \dim B_{3'}=9,
\]
并且
\[
A_2\sim_{\QQ} B_2^{2},\qquad
A_3\sim_{\QQ} B_3^{3},\qquad
A_{3'}\sim_{\QQ} B_{3'}^{3}.
\]
从而
\[
J\sim_{\QQ} A_{\mathrm{sgn}}\times B_2^{2}\times B_3^{3}\times B_{3'}^{3}.
\]
\end{theorem}
\begin{proof}
对 \(\rho\in\{V_2,V_3,V_3'\}\)，记 \(A_\rho\) 为 \(J\) 的 \(\rho\)-等型部分。由既有的 \(S_4\)-等型分解理论，\(\End_{S_4}(\rho)\cong \QQ\)，故群代数 \(\QQ[S_4]\) 在 \(A_\rho\) 上的像同构于矩阵代数 \(M_{\dim \rho}(\QQ)\)。取其中任一秩一幂等元 \(f_\rho\)，定义
\[
B_\rho:=f_\rho A_\rho.
\]
于是标准矩阵代数论证给出
\[
A_\rho\sim_{\QQ} B_\rho^{\dim \rho}.
\]
对 \(\rho=V_2,V_3,V_3'\) 分别应用，即得
\[
A_2\sim_{\QQ}B_2^{2},\qquad
A_3\sim_{\QQ}B_3^{3},\qquad
A_{3'}\sim_{\QQ}B_{3'}^{3}.
\]
又由全纯微分表示分解，
\[
\dim A_2=4\cdot 2=8,\qquad
\dim A_3=3\cdot 3=9,\qquad
\dim A_{3'}=9\cdot 3=27,
\]
故
\[
\dim B_2=4,\qquad \dim B_3=3,\qquad \dim B_{3'}=9.
\]
代回即可得到 \(J\) 的根块化分解。证毕。
\end{proof}

\begin{theorem}[任意子群商的根块公式]\label{thm:conclusion-s4-arbitrary-subgroup-quotient-root-block-formula}
对任意子群 \(H\le S_4\)，定义
\[
d_{\mathrm{sgn}}(H):=\dim(\varepsilon^H),\qquad
d_{\rho_2}(H):=\dim(V_2^H),\qquad
d_{\rho_3}(H):=\dim(V_3^H),\qquad
d_{\rho_3'}(H):=\dim((V_3')^H),
\]
亦即
\[
d_\rho(H)=\frac{1}{|H|}\sum_{h\in H}\chi_\rho(h).
\]
则有 \(\QQ\)-同源分解
\[
\Jac(X/H)\sim_{\QQ}
A_{\mathrm{sgn}}^{\,d_{\mathrm{sgn}}(H)}
\times
B_2^{\,d_{\rho_2}(H)}
\times
B_3^{\,d_{\rho_3}(H)}
\times
B_{3'}^{\,d_{\rho_3'}(H)}.
\]
因此
\[
g(X/H)=
5\,d_{\mathrm{sgn}}(H)+
4\,d_{\rho_2}(H)+
3\,d_{\rho_3}(H)+
9\,d_{\rho_3'}(H).
\]
\end{theorem}
\begin{proof}
令
\[
e_H:=\frac{1}{|H|}\sum_{h\in H}h\in \QQ[S_4].
\]
按 pullback--norm 的标准机制，\(e_H\) 在 \(\QQ\)-同源范畴中切出 \(H\)-不变部分，故
\[
e_HJ\sim_{\QQ}\Jac(X/H).
\]
另一方面，由定理 \ref{thm:conclusion-s4-jacobian-root-block-realization}，
\[
J\sim_{\QQ}A_{\mathrm{sgn}}\times B_2^{2}\times B_3^{3}\times B_{3'}^{3}.
\]
在 \(\rho\)-等型部分上，\(e_H\) 的秩恰为 \(\dim \rho^H=d_\rho(H)\)。因此在根块分解中，\(e_H\) 分别切出
\[
A_{\mathrm{sgn}}^{\,d_{\mathrm{sgn}}(H)},\qquad
B_2^{\,d_{\rho_2}(H)},\qquad
B_3^{\,d_{\rho_3}(H)},\qquad
B_{3'}^{\,d_{\rho_3'}(H)}.
\]
四个等型部分相乘即得所述同源公式。亏格公式由
\[
\dim A_{\mathrm{sgn}}=5,\qquad \dim B_2=4,\qquad \dim B_3=3,\qquad \dim B_{3'}=9
\]
直接相加得到。证毕。
\end{proof}

\begin{theorem}[标准子群商的完整 atlas]\label{thm:conclusion-s4-standard-subgroup-rootblock-atlas}
取标准子群代表
\[
C_2^{\mathrm t}:=\langle(13)\rangle,\qquad
C_2^{\mathrm d}:=\langle(13)(24)\rangle,\qquad
C_3:=\langle(123)\rangle,\qquad
C_4:=\langle(1234)\rangle,
\]
\[
V_4^{\mathrm n}:=\langle(12)(34),(13)(24)\rangle,\qquad
V_4^{\mathrm m}:=\langle(13),(24)\rangle,
\]
\[
S_3:=\langle(12),(123)\rangle,\qquad
D_8:=\langle(1234),(13)\rangle,\qquad
A_4.
\]
则有下列精确同源分解与亏格表：
\[
\begin{array}{c|c|c}
H & \Jac(X/H)\sim_{\QQ} & g(X/H)\\
\hline
C_2^{\mathrm t} & B_2\times B_3^{2}\times B_{3'} & 19\\
C_2^{\mathrm d} & A_{\mathrm{sgn}}\times B_2^{2}\times B_3\times B_{3'} & 25\\
C_3 & A_{\mathrm{sgn}}\times B_3\times B_{3'} & 17\\
C_4 & B_2\times B_{3'} & 13\\
V_4^{\mathrm n} & A_{\mathrm{sgn}}\times B_2^{2} & 13\\
V_4^{\mathrm m} & B_2\times B_3 & 7\\
S_3 & B_3 & 3\\
D_8 & B_2 & 4\\
A_4 & A_{\mathrm{sgn}} & 5
\end{array}
\]
\end{theorem}
\begin{proof}
在共轭类
\[
1,\ (12),\ (12)(34),\ (123),\ (1234)
\]
上的角色值为
\[
\chi_{\varepsilon}=(1,-1,1,1,-1),
\]
\[
\chi_{V_2}=(2,0,2,-1,0),
\]
\[
\chi_{V_3}=(3,1,-1,0,-1),
\]
\[
\chi_{V_3'}=(3,-1,-1,0,1).
\]
将这些角色值在各子群上平均，即得相应的不变维数 \(d_\rho(H)\)。例如
\[
d_{\rho_2}(D_8)=\frac18\bigl(2+2\cdot 0+3\cdot 2+2\cdot 0\bigr)=1,
\]
而其余三个不可约表示在 \(D_8\) 上的不变维数皆为零，故
\[
\Jac(X/D_8)\sim_{\QQ}B_2.
\]
其余各行完全同理。将所得 \(d_\rho(H)\) 代入定理 \ref{thm:conclusion-s4-arbitrary-subgroup-quotient-root-block-formula}，即得全表；亏格为各因子维数之和。证毕。
\end{proof}

\begin{proposition}[\texorpdfstring{$S_4$}{S4} 作用的固定点 census]\label{prop:conclusion-s4-fixedpoint-census-repackaged}
对共轭类代表有
\[
\#\operatorname{Fix}_X((12))=24,\qquad
\#\operatorname{Fix}_X((12)(34))=0,\qquad
\#\operatorname{Fix}_X((123))=0,\qquad
\#\operatorname{Fix}_X((1234))=0.
\]
等价地，仅有换位在 \(X\) 上出现几何固定点；全部双换位、\(3\)-循环与 \(4\)-循环都自由作用。
\end{proposition}
\begin{proof}
这正是前述固定点 census 的直接重述。证毕。
\end{proof}

\begin{theorem}[商曲线亏格的换位计数律]\label{thm:conclusion-s4-quotient-genus-transposition-count-law}
\leanverified{paper\_conclusion\_s4\_quotient\_genus\_transposition\_count\_law}
对任意子群 \(H\le S_4\)，记
\[
t(H):=\#\bigl(H\cap\{\text{换位}\}\bigr).
\]
则
\[
g(X/H)=1+\frac{48-12\,t(H)}{|H|}.
\]
特别地，
\[
X\to X/H\ \text{为 \'{e}tale}
\iff
t(H)=0.
\]
\end{theorem}
\begin{proof}
由命题 \ref{prop:conclusion-s4-fixedpoint-census-repackaged}，只有换位可能产生固定点。若两不同换位 \(\tau_1\neq\tau_2\) 在某点 \(x\in X\) 上同时固定，则 \(\tau_1\tau_2\) 亦固定 \(x\)。而 \(\tau_1\tau_2\) 非单位时只能是 \(3\)-循环或双换位，这与命题 \ref{prop:conclusion-s4-fixedpoint-census-repackaged} 矛盾。故不同换位的固定点集两两不交。

每个换位恰固定 \(24\) 个点，因此对任意 \(H\)，总分歧贡献为
\[
24\,t(H).
\]
对商映射 \(X\to X/H\) 应用 Riemann--Hurwitz，
\[
2g(X)-2
=
|H|\bigl(2g(X/H)-2\bigr)+24\,t(H).
\]
由于 \(g(X)=49\)，即
\[
96=|H|\bigl(2g(X/H)-2\bigr)+24\,t(H).
\]
整理即得
\[
g(X/H)=1+\frac{48-12\,t(H)}{|H|}.
\]
末句显然，因为分歧为空当且仅当总分歧贡献为零，而这等价于 \(t(H)=0\)。证毕。
\end{proof}

\begin{corollary}[proper quotient 的 genus gap 与唯一极小实现]\label{cor:conclusion-s4-proper-quotient-minimal-genus-realization}
在全部真子群商 \(X/H\) 中，最小的三个亏格恰为
\[
3,\qquad 4,\qquad 5,
\]
并且分别唯一地由
\[
H\cong S_3,\qquad H\cong D_8,\qquad H\cong A_4
\]
实现。相应地，
\[
\Jac(X/S_3)\sim_{\QQ} B_3,\qquad
\Jac(X/D_8)\sim_{\QQ} B_2,\qquad
\Jac(X/A_4)\sim_{\QQ} A_{\mathrm{sgn}}.
\]
因此，三块最小可见根块 \(B_3,B_2,A_{\mathrm{sgn}}\) 分别由亏格 \(3,4,5\) 的商曲线唯一实现。
\end{corollary}
\begin{proof}
由定理 \ref{thm:conclusion-s4-standard-subgroup-rootblock-atlas}，真子群商的亏格最小值依次为
\[
g(X/S_3)=3,\qquad g(X/D_8)=4,\qquad g(X/A_4)=5,
\]
其余标准子群商亏格均严格大于 \(5\)。相应的 Jacobian 分解亦由同一定理直接给出。证毕。
\end{proof}

\begin{theorem}[两条 Prym 梯：\texorpdfstring{$B_3$}{B3} 的 \'{e}tale realization 与 \texorpdfstring{$B_{3'}$}{B3'} 的 ramified realization]\label{thm:conclusion-s4-two-prym-ladders}
设
\[
Y:=X/V_4^{\mathrm m},\qquad Z:=X/D_8,\qquad W:=X/C_4.
\]
则存在两条自然二次覆盖
\[
Y\longrightarrow Z,\qquad W\longrightarrow Z,
\]
并且：
\begin{enumerate}
  \item \(Y\to Z\) 为 \'{e}tale 双覆盖，且
  \[
  \operatorname{Prym}(Y/Z)\sim_{\QQ} B_3,\qquad \dim \operatorname{Prym}(Y/Z)=3;
  \]
  \item \(W\to Z\) 为分歧点数 \(12\) 的双覆盖，且
  \[
  \operatorname{Prym}(W/Z)\sim_{\QQ} B_{3'},\qquad \dim \operatorname{Prym}(W/Z)=9.
  \]
\end{enumerate}
\end{theorem}
\begin{proof}
由于
\[
V_4^{\mathrm m}\triangleleft D_8,\qquad C_4\triangleleft D_8,
\]
且两者指数都为 \(2\)，故相应二次覆盖存在。

先看 \(Y\to Z\)。由定理 \ref{thm:conclusion-s4-standard-subgroup-rootblock-atlas}，
\[
g(Y)=7,\qquad g(Z)=4.
\]
Riemann--Hurwitz 给出
\[
2g(Y)-2=2(2g(Z)-2)+r.
\]
代入得到
\[
12=12+r,
\]
故 \(r=0\)，因此 \(Y\to Z\) \'{e}tale。又由同一定理，
\[
\Jac(Y)\sim_{\QQ} B_2\times B_3,\qquad \Jac(Z)\sim_{\QQ} B_2.
\]
范数映射的连通核同源于互补因子 \(B_3\)，故
\[
\operatorname{Prym}(Y/Z)\sim_{\QQ} B_3.
\]
其维数为 \(g(Y)-g(Z)=3\)。

再看 \(W\to Z\)。同样由定理 \ref{thm:conclusion-s4-standard-subgroup-rootblock-atlas}，
\[
g(W)=13,\qquad g(Z)=4.
\]
Riemann--Hurwitz 给出
\[
2g(W)-2=2(2g(Z)-2)+r,
\]
即
\[
24=12+r,
\]
从而 \(r=12\)。另一方面，既有九维 Prym 结论给出
\[
\Jac(W)\sim_{\QQ} B_2\times B_{3'},\qquad \Jac(Z)\sim_{\QQ} B_2,
\]
于是
\[
\operatorname{Prym}(W/Z)\sim_{\QQ} B_{3'}.
\]
其维数即 \(g(W)-g(Z)=9\)。证毕。
\end{proof}

\begin{theorem}[低亏格 exact isogeny identities]\label{thm:conclusion-s4-low-genus-exact-isogeny-identities}
记
\[
P_9:=\operatorname{Prym}(X/C_4\to X/D_8)\sim_{\QQ} B_{3'}.
\]
则存在如下精确 \(\QQ\)-同源恒等式：
\[
\Jac(X/V_4^{\mathrm m})
\sim_{\QQ}
\Jac(X/D_8)\times \Jac(X/S_3),
\]
\[
\Jac(X/C_4)
\sim_{\QQ}
\Jac(X/D_8)\times P_9,
\]
\[
\Jac(X/C_3)
\sim_{\QQ}
\Jac(X/A_4)\times \Jac(X/S_3)\times P_9,
\]
\[
\Jac(X/C_2^{\mathrm t})
\sim_{\QQ}
\Jac(X/D_8)\times \Jac(X/S_3)^{2}\times P_9.
\]
\end{theorem}
\begin{proof}
由推论 \ref{cor:conclusion-s4-proper-quotient-minimal-genus-realization} 与定理 \ref{thm:conclusion-s4-two-prym-ladders}，
\[
\Jac(X/D_8)\sim_{\QQ} B_2,\qquad
\Jac(X/S_3)\sim_{\QQ} B_3,\qquad
\Jac(X/A_4)\sim_{\QQ} A_{\mathrm{sgn}},\qquad
P_9\sim_{\QQ} B_{3'}.
\]
将这些代入定理 \ref{thm:conclusion-s4-standard-subgroup-rootblock-atlas} 中
\[
\Jac(X/V_4^{\mathrm m})\sim_{\QQ} B_2\times B_3,
\]
\[
\Jac(X/C_4)\sim_{\QQ} B_2\times B_{3'},
\]
\[
\Jac(X/C_3)\sim_{\QQ} A_{\mathrm{sgn}}\times B_3\times B_{3'},
\]
\[
\Jac(X/C_2^{\mathrm t})\sim_{\QQ} B_2\times B_3^{2}\times B_{3'},
\]
即得所列精确同源恒等式。证毕。
\end{proof}

\begin{theorem}[属 \texorpdfstring{$49$}{49} 的 \texorpdfstring{$S_4$}{S4} 闭包 Jacobian 的低亏格全局压缩]\label{thm:conclusion-s4-global-jacobian-low-genus-compression}
记
\[
P_9:=\operatorname{Prym}(X/C_4\to X/D_8).
\]
则整个 Jacobian 满足
\[
J
\sim_{\QQ}
\Jac(X/A_4)\times
\Jac(X/D_8)^{2}\times
\Jac(X/S_3)^{3}\times
P_9^{3}.
\]
等价地，
\[
J\sim_{\QQ}
A_{\mathrm{sgn}}\times B_2^{2}\times B_3^{3}\times B_{3'}^{3}.
\]
因而 \(49\) 维整体 Jacobian 已完全压缩为亏格 \(5,4,3\) 的三个商曲线与一条 \(13\to4\) 的二次 Prym 接口所承载的数据。
\end{theorem}
\begin{proof}
定理 \ref{thm:conclusion-s4-jacobian-root-block-realization} 给出
\[
J\sim_{\QQ}
A_{\mathrm{sgn}}\times B_2^{2}\times B_3^{3}\times B_{3'}^{3}.
\]
又由推论 \ref{cor:conclusion-s4-proper-quotient-minimal-genus-realization}，
\[
A_{\mathrm{sgn}}\sim_{\QQ}\Jac(X/A_4),\qquad
B_2\sim_{\QQ}\Jac(X/D_8),\qquad
B_3\sim_{\QQ}\Jac(X/S_3),
\]
而由定理 \ref{thm:conclusion-s4-two-prym-ladders}，
\[
B_{3'}\sim_{\QQ}P_9.
\]
逐项代回即得全局压缩式。证毕。
\end{proof}

由此，\(S_4\)-正则闭包曲线的 Jacobian 不再只是由一条 \(49\) 维整体对象承载。全纯微分的 Chevalley--Weil 等型分解先经根块化被压为 \(A_{\mathrm{sgn}},B_2,B_3,B_{3'}\) 四类基本因子；任意子群商的 Jacobian 再由不变维数公式一键读出；固定点 census 与换位计数律把全部分歧数据压回子群内部的换位计数；而低亏格端点恰由 \(X/S_3\)、\(X/D_8\)、\(X/A_4\) 三条商曲线与 \(X/C_4\to X/D_8\) 的九维 Prym 梯完全实现。于是，整个 genus \(49\) 的 \(S_4\)-闭包 Jacobian 在 \(\QQ\)-同源范畴中已被低亏格 quotient atlas 与单条 Prym 接口完全压缩。

\endinput
